\clearpage
\begin{trapbox}
\paragraph{误区一：把旧余弦错式当作低精度的正确公式}
\begin{center}
\begin{tikzpicture}[font=\small]
\draw[black!30,rounded corners=2pt] (0,0) rectangle (14.4,8.2);
\node[anchor=west,font=\small\bfseries] at (0.5,7.72) {余弦错式对照：旧错式破坏了偶对称};
\begin{axis}[width=11.2cm,height=6.1cm,at={(2.05cm,1.25cm)},anchor=south west,
axis lines=left,xlabel={$x$},tick label style={font=\small},label style={font=\small},
legend style={font=\small,draw=none,fill=white,at={(0.03,0.97)},anchor=north west},
legend cell align=left,unbounded coords=jump,xmin=-0.8,xmax=0.8,ymin=0.55,ymax=1.65,ylabel={$y$},legend style={at={(0.97,0.97)},anchor=north east}]
\addplot[blue,very thick] table[x=x,y=cos] {assets/tikz/F003_samples.dat};
\addlegendentry{$\cos x$}
\addplot[orange!85!black,thick,dashed] table[x=x,y=rc] {assets/tikz/F003_samples.dat};
\addlegendentry{正确 $R_C(x)$}
\addplot[red!70!black,thick,densely dotted,domain=-0.8:0.8,samples=100]{12/(12+5*x)};\addlegendentry{错误 $12/(12+5x)$}
\end{axis}
\end{tikzpicture}
\end{center}

错误分式 $W(x)=12/(12+5x)$ 不具备偶对称，且 $W'(0)=-5/12\ne0$，而 $\cos'(0)=0$。
所以它不是本条的余弦帕德公式，不能以“精度低一些”为由保留。
正确替代是 $R_C=(12-5x^2)/(12+x^2)$，它保持偶对称并匹配至四次项。

\newpage
\paragraph{误区二：只排除分式分母为零，就认为原函数也有定义}
\begin{center}
\begin{tikzpicture}[font=\small]
\draw[black!30,rounded corners=2pt] (0,0) rectangle (14.4,8.2);
\node[anchor=west,font=\small\bfseries] at (0.5,7.72) {正切的真实极点与分式极点不重合};
\begin{axis}[width=11.2cm,height=6.1cm,at={(2.05cm,1.25cm)},anchor=south west,
axis lines=left,xlabel={$x$},tick label style={font=\small},label style={font=\small},
legend style={font=\small,draw=none,fill=white,at={(0.03,0.97)},anchor=north west},
legend cell align=left,unbounded coords=jump,xmin=1.1,xmax=1.9,ymin=-35,ymax=35,ylabel={$y$},legend style={at={(0.03,0.05)},anchor=south west}]

\addplot[blue,very thick,domain=1.1:1.543,samples=100]{tan(deg(x))};\addlegendentry{$\tan x$}
\addplot[blue,very thick,forget plot,domain=1.599:1.9,samples=100]{tan(deg(x))};
\addplot[orange!85!black,thick,dashed,domain=1.1:1.70,samples=100]{3*x/(3-x*x)};\addlegendentry{$R_T(x)$}
\addplot[orange!85!black,thick,dashed,forget plot,domain=1.76:1.9,samples=100]{3*x/(3-x*x)};
\draw[blue,densely dotted] (axis cs:1.570796327,-35)--(axis cs:1.570796327,24);
\draw[orange!85!black,densely dotted] (axis cs:1.732050808,-35)--(axis cs:1.732050808,24);
\node[blue] at (axis cs:1.570796327,29) {$\pi/2$};
\node[orange!85!black] at (axis cs:1.732050808,29) {$\sqrt3$};

\end{axis}
\end{tikzpicture}
\end{center}

$\tan x$ 在 $\pi/2$ 无定义，但 $R_T$ 在该处仍有定义；$R_T$ 的正极点却是 $\sqrt3$。
正确做法是\textbf{同时}检查原函数和分式的定义域，并保持在适合局部估算的范围。不同极点位置也直接否定了全局恒等的想法。

\newpage
\paragraph{误区三：分式处处有定义，就能模拟全局变化}
\begin{center}
\begin{tikzpicture}[font=\small]
\draw[black!30,rounded corners=2pt] (0,0) rectangle (14.4,8.2);
\node[anchor=west,font=\small\bfseries] at (0.5,7.72) {远离零点：分式不能替代指数的整体增长};
\begin{axis}[width=11.2cm,height=6.1cm,at={(2.05cm,1.25cm)},anchor=south west,
axis lines=left,xlabel={$x$},tick label style={font=\small},label style={font=\small},
legend style={font=\small,draw=none,fill=white,at={(0.03,0.97)},anchor=north west},
legend cell align=left,unbounded coords=jump,xmin=0,xmax=6,ymin=0,ymax=450,ylabel={$y$}]

\addplot[blue,very thick,domain=0:6,samples=150]{exp(x)};\addlegendentry{$e^x$}
\addplot[orange!85!black,thick,dashed,domain=0:6,samples=150]{(x*x+6*x+12)/(x*x-6*x+12)};\addlegendentry{$R_E(x)$}

\end{axis}
\end{tikzpicture}
\end{center}

$R_E$ 的分母始终为正，然而其导数
\[
R_E'(x)=\frac{12(12-x^2)}{(x^2-6x+12)^2}
\]
在 $x>\sqrt{12}$ 时为负，指数函数却始终递增。
正确替代：先指定区间与容许误差，再用误差界检验；不把“$|x|<1$”当作保证任意精度的万能条件。

\paragraph{误区四：角度数字直接代入，或把近似当作恒等}
$30^\circ$ 应代入 $x=\pi/6$，而不是 $30$；单位圆图已显示其弧度对应关系。
近似式用于证明等式或不等式前，必须另证误差为零或确定误差方向，不能把 $\approx$ 换成 $=$。
\end{trapbox}
